\subsection{DeformNeXt-10k construction}
\label{sec:data-construction}
We construct paired shape edits from the textured assets and hierarchical
part annotations of PartNeXt~\cite{wang2025partnext}. A canonical source
$M_0=(V_0,F,\mathcal A_0)$ defines vertex indices, ordered faces, and the
UV/material/corner-attribute authority. Each accepted variant is
$M_{1,k}=(V_{1,k},F,\mathcal A_0)$: positions change while correspondence
and face connectivity/order remain fixed. Shared face indices and barycentric
coordinates therefore identify corresponding source and edited surface
points without nearest-neighbor reassignment.

\paragraph{Donor-guided proposals.}
We select donors from the same category using frozen semantic-role and
structural-compatibility heuristics. Semantic roles can contain multiple
disconnected pieces; matching individual leaf meshes or requiring equal
part counts would unnecessarily restrict shape diversity. Source/donor
identities and configuration hashes are recorded before each run. Eight
original-appearance donor renders personalize an SDXL prior with
GeoStyle~\cite{choi2026geostyle}. The donor is a morphology prior, rather
than target geometry or an input to the retargeting model.

GeoStyle proposes face-Jacobian changes under image-space score distillation.
Our production path adds a category prompt and source assembly constraints
to the reconstruction: selected primary surfaces deform, complementary
physical components follow rigid motions, and source-derived attachment
proxies couple their movement. These contacts are geometric estimates,
not annotated physical joints. The main production configuration uses
800 donor-specific LoRA updates and 300 deformation iterations; donor LoRA
and source preprocessing are reused across pairs. This shortened configuration
was adopted for throughput and is not demonstrated equivalent to the
official 3,600-iteration configuration. Supplementary details give the
other changes and retained early configurations.

\paragraph{Final-pair selection.}
The category-prompt/assembly proposal (PA) and an optional 120-iteration
Jacobian, crease-aware normal, and neighboring-displacement refinement
(balanced JN-L) are separate candidates. Each is assessed anew against
$M_0$. Agreement with the initial proposal or the donor is not an admission
criterion. We check canonical correspondence, finite geometry and
catastrophic failures, meaningful deformation, paired RGB usability,
native-representation survival, and actual training-loader compatibility.
Multi-view model-assisted screening identifies obvious spikes, detached
components and assembly collapse; it is not independent human validation.
Small intersections, gaps and folds are recorded as diagnostics rather
than universally rejected. The dataset is therefore not certified
intersection-free. Both meshes retain one source-defined 0.48 frame and
are rendered with the source materials and matched cameras; edited meshes
are not independently normalized. All approved distinct variants are
retained, with near-identical final geometries merged.

\begin{figure*}[t]
\centering
\includegraphics[width=\textwidth]{figures/deformnext-construction.pdf}
\caption{DeformNeXt construction and approved source/edited pairs. Donor
images guide morphology; source assembly constrains reconstruction, and each
final candidate is screened independently. Examples show archived RGB
(camera 8) and clay (camera 14), with a common foreground-union display crop
and scale for M0/M1 at each camera. Four categories are drawn from a pre-existing
large-edit illustration selection whose membership was checked against the
final snapshot. They illustrate visible changes, not typical quality,
generator yield or a 300-versus-3,600-step comparison.}
\label{fig:deformnext-construction}
\end{figure*}

\paragraph{Scope and splits.}
The final snapshot contains 10,165 distinct deformations from 3,483 source
UIDs across 50 categories, averaging 2.92 variants per source. This is a
deformation count, not 10,000 independent source objects or an estimate of
yield over all PartNeXt assets. After construction, a fixed seed assigns
approximately 90/5/5 of the examples to train/validation/test while keeping
every source UID and source geometry cluster in one split
(Table~\ref{tab:deformnext-splits}). Donor disjointness is not enforced;
the held-out evaluation measures retargeting on admitted edits from this
construction process, not generalization to an unseen deformation generator.

\begin{table}[t]
\centering
\small
\caption{Frozen DeformNeXt splits. Variants of one source and geometry
cluster stay together; category proportions are not forced to be uniform.}
\label{tab:deformnext-splits}
\begin{tabular}{lrrrr}
\toprule
Split & Edits & UIDs & Clusters & Classes \\
\midrule
Train & 9,149 & 3,184 & 1,790 & 50 \\
Validation & 508 & 153 & 132 & 42 \\
Test & 508 & 146 & 134 & 42 \\
\midrule
Total & 10,165 & 3,483 & 2,056 & 50 \\
\bottomrule
\end{tabular}
\end{table}
